\section{Gait pattern}

\hfill \break
The gait pattern, also known as the gait cycle, refers to the sequence of movements that occur when a person walks or runs. The gait cycle can be broken down into two main phases: the stance phase and the swing phase.\\

The stance phase is the portion of the gait cycle in which the foot is in contact with the ground. During this phase, the body's weight is transferred from the heel to the toe, and the knee and hip joints work to provide stability and propulsion. There are three sub-phases within the stance phase: heel strike, mid-stance, and toe-off.\\
The swing phase is the portion of the gait cycle in which the foot is off the ground. During this phase, the hip, knee, and ankle joints work together to bring the foot forward and prepare it for the next heel strike.\\

Normal gait pattern is considered symmetric, smooth and efficient. Any abnormalities or variations in the gait pattern can be an indication of certain medical conditions or injuries. Common examples include limping, which is a variation in the gait pattern that can be caused by pain or weakness in the leg, or a shuffling gait which can be caused by conditions like Parkinson's disease.\\

Long term goal of this project is to detect analyse a person´s gait pattern and detect if the person is healthy. however, first, detecting the type of motion will be investigated.
To do so Walking, running, and jumping will be analysed and a record of poeple performing these motion will be used as a data set for our Tool. 
These three forms of human locomotion, have some key differences in terms of the mechanics of the movements involved.\\

Running is a fast and vigorous form of movement, and is characterized by both feet leaving the ground during the gait cycle. During running, the foot typically strikes the ground with the heel and then quickly rolls forward to push off with the toe, it has a greater stride length compared to Walking, and a segnificant thigh angle.\\

Walking is characterized by one foot being in contact with the ground at all times. During walking, the foot typically strikes the ground with the heel first and then rolls forward to push off with the toe. Walking is a relatively slow and steady form of movement, it has a greater stride width than Running, and a smaller thigh angle.\\

Jumping is a sudden, explosive movement where both feet leave the ground at the same time and is most commonly used for vertical displacement, the thigh angle is . Jumping is a complex and highly coordinated movement, which can be performed in many ways such as for example for athletic events like high jump, long jump or pole vault. \\

All three forms of movement have their own peculiarities that can be measured with an IMU: running and walking are distinguished by speed and thigh angle, while jumping is characterized by the direction of movement.\\
In walking, the thigh angle typically decreases during the stance phase and increases during the swing phase, while in running, the thigh angle typically increases during the stance phase and decreases during the swing phase. Jumping involves a larger thigh angle excursion than both walking and running.

Additionally, the accelerometer present on our board can be used to detect and classify walking, running, and jumping. Based on the acceleration of the body segment, the algorithm can distinguish the type of motion.\\

However, it is important to note that while the thigh angle can provide some information about a person's gait pattern, it is not the only factor that needs to be considered. Other factors such as the foot trajectory, ankle angle and the knee angle also play an important role in identifying walking, running and jumping. Also, It is also important to consider that this is a very complex process and there is much research ongoing to improve the accuracy of gait detection and classification. Furthermore, there are many factors that can influence the thigh angle such as shoes, surfaces, and age among others.


%https://slideplayer.com/slide/12956873/


